\chapter{پیاده سازی و تحلیل نتایج}
\label{chap:implementation}

\section{داده عمومی و قابلیت بازتولید}
آزمایش با ۳۲۸۹ رکورد ساعتی، از ساعت ۱۷ روز ۱۱ ژانویه تا ساعت ۱۷ روز ۲۷ مه ۲۰۱۶ به وقت محلی \lr{Stambruges}، انجام شده است. این شمار از گام‌های ساعتی کامل همان بازه به دست آمده و با جدول~\ref{tab:data-origin} سازگار است. همان جدول نشان می‌دهد این مجموعه «داده عمومی ترکیبی» است، نه اندازه‌گیری یک ریزشبکه در محل. پارامترهای بایگانی‌شده آزمایش در پیوست~\ref{app:reproducibility} آمده‌اند.

\begin{table}[H]
  \centering
  \caption{منشأ مؤلفه‌های مطالعه موردی}
  \label{tab:data-origin}
  \begin{tabular}{>{\centering\arraybackslash}p{0.18\textwidth}p{0.40\textwidth}p{0.30\textwidth}}
    \toprule
    \textbf{مؤلفه} & \textbf{منشأ و تبدیل} & \textbf{محدودیت} \\
    \midrule
    بازه و شمار رکورد &
    ۱۱ ژانویه ۲۰۱۶، ساعت ۱۷ تا ۲۷ مه ۲۰۱۶، ساعت ۱۷؛ ۳۲۸۹ رکورد ساعتی &
    بازه پردازش‌شده است و نباید با بازه کامل مجموعه اصلی یکسان فرض شود \\
    بار &
    مدل بار مسکونی روی تقویم کانداندو با بار زمینه $0{.}20\,\mathrm{kW}$؛ اگر فایل \lr{UCI} در دسترس باشد میانگین ساعتی وسایل و روشنایی جایگزین می‌شود &
    بار کامل یک محله یا ریزشبکه نیست \\
    فتوولتائیک &
    پوش ارتفاع خورشید در \lr{Stambruges}؛ توان نامی $1\,\mathrm{kW}$ و نسبت عملکرد $0{.}82$؛ اگر خدمت \lr{NASA POWER} در دسترس باشد تابش روزانه آن پخش می‌شود &
    توان فتوولتائیک ساعتی اندازه‌گیری‌شده در همان خانه نیست \\
    تعرفه &
    تعرفه زمان‌استفاده بایگانی‌شده در پیوست~\ref{app:reproducibility} &
    قیمت بازار واقعی یا قبض مشترک نیست \\
    باتری &
    مدل انرژی با ظرفیت $2\,\mathrm{kWh}$ &
    دما، پیری و محدودیت مبدل ندارد \\
    \bottomrule
  \end{tabular}
\end{table}

تقویم و محل مطالعه از مجموعه عمومی مصرف وسایل یک خانه کم‌مصرف در \lr{Stambruges} گرفته شده است \pcite{candanedo2017energy,candanedo2017uci}. اسکریپت \lr{code/prepare\_public\_data.py} اگر فایل \lr{UCI} و خدمت \lr{NASA POWER} در دسترس باشند بار ساعتی و تابش روزانه را از آن‌ها می‌سازد؛ در غیر این صورت همان بازه و مختصات را با مدل بار مسکونی و پوش ارتفاع خورشید پر می‌کند. مجموعه حاصل مطالعه موردی ترکیبی و بازتولیدپذیر است و نباید «آزمون داده میدانی ریزشبکه» نامیده شود.

کد شبیه‌سازی در شاخه \lr{code} است. اسکریپت \lr{prepare\_public\_data.py} سری ساعتی را می‌سازد، \lr{test\_physics.py} لایه فیزیکی را در حالت‌های حدی می‌آزماید و \lr{simulate.py} مرجع برنامه‌ریزی پویا، آموزش پنج بذر و شکل‌ها را تولید می‌کند.

\section{تنظیمات اجرا}
از ۳۲۸۹ ساعت، ۱۹۷۳ ساعت آموزش، ۶۵۷ ساعت اعتبارسنجی و ۶۵۹ ساعت پایانی آزمون است. افراز کاملاً زمانی است تا الگوی آینده به گذشته نشت نکند. هر دو شبکه دو لایه پنهان ۳۲ واحدی دارند. شبکه داده‌محور چهار خروجی آزاد و شبکه فیزیک‌تعبیه‌شده یک خروجی فرمان دارد. آموزش حداکثر ۴۵۰ دوره با دسته ۶۴تایی است و بهترین دوره هر بذر فقط با خطای اعتبارسنجی انتخاب می‌شود.

مرجع برنامه‌ریزی پویا وضعیت شارژ را با گام $0{.}01$ در بازه $[0{.}20,0{.}90]$ و فرمان باتری را با ۴۱ مقدار در بازه $[-1,1]\,\mathrm{kW}$ جست‌وجو می‌کند. پس مرجع بهینه پیوسته دقیق نیست؛ اما یک مبنای مقید، صریح و مستقل از شبکه عصبی است.

\section{نتایج دقت و قیود}
\textbf{حالت بدون باتری} یعنی شارژ و دشارژ صفر است و همه عدم‌توازن با شبکه بسته می‌شود: $P_{\mathrm{c}}=P_{\mathrm{d}}=0$ و $P_{\mathrm{G}}=P_{\mathrm{L}}-P_{\mathrm{PV}}$. وضعیت شارژ در مقدار آغازین افق آزمون ثابت می‌ماند. در این ردیف برابری $\mathrm{MAE}(P_{\mathrm{G}})$ و $\mathrm{MAE}(u)$ از $P_{\mathrm{G}}-P_{\mathrm{G}}^{\star}=-u^{\star}$ با انحصار فرمان مرجع می‌آید.

جدول~\ref{tab:accuracy} خطای تقریب نسبت به برنامه مرجع را نشان می‌دهد. ردیف مرجع نسبت به خودش صفر است و فقط مقیاس جدول را مشخص می‌کند. \lr{MSE} مطابق~\eqref{eq:mse-comp} میانگین چهار مؤلفه است و برای مقایسه عددی است؛ \lr{MAE} با واحد جداگانه خواناتر است. مدل داده‌محور آزاد \lr{MSE} پایین‌تری دارد، اما این به معنی قابل اجرا بودن فرمان نیست.

\begin{table}[H]
  \centering
  \caption{میانگین و انحراف معیار پنج بذر روی افق آزمون}
  \label{tab:accuracy}
  \begin{tabular}{lcccc}
    \toprule
    مدل & $\mathrm{MSE}_{\mathrm{comp}}$ & $\mathrm{MAE}(P_{\mathrm{G}})$ (کیلووات) & $\mathrm{MAE}(u)$ (کیلووات) & $\mathrm{MAE}(s)$ \\
    \midrule
    مرجع پویا & $0$ & $0$ & $0$ & $0$ \\
    بدون باتری & $0{.}0444$ & $0{.}099$ & $0{.}099$ & $0{.}263$ \\
    داده‌محور خام & $0{.}0133\pm 0{.}0008$ & $0{.}079\pm 0{.}003$ & $0{.}078\pm 0{.}004$ & $0{.}093\pm 0{.}009$ \\
    فیزیک‌تعبیه‌شده & $0{.}0145\pm 0{.}0016$ & $0{.}077\pm 0{.}008$ & $0{.}077\pm 0{.}008$ & $0{.}097\pm 0{.}017$ \\
    \bottomrule
  \end{tabular}
\end{table}

جدول~\ref{tab:feasibility} نشان می‌دهد چرا \lr{MSE} به‌تنهایی کافی نیست. در خروجی خام مدل داده‌محور حدود $11{.}4$ درصد مقادیر وضعیت شارژ نامعتبرند و به‌ترتیب $41{.}9$ درصد و $31{.}0$ درصد مقادیر شارژ و دشارژ منفی‌اند. در حدود $35{.}4$ درصد گام‌ها هر دو توان هم‌زمان مثبت‌اند. لایه فیزیک‌تعبیه‌شده همه این نرخ‌ها را در محاسبه شناور به صفر می‌رساند، چون قیدها در خود خروجی‌اند. این صفر بودن کشف آماری نیست؛ خاصیت طراحی است.

\begin{table}[H]
  \centering
  \caption{شاخص‌های امکان‌پذیری روی مجموعه آزمون}
  \label{tab:feasibility}
  \begin{tabular}{lccccc}
    \toprule
    مدل & $\overline{|r_{\mathrm{G}}|}$ (کیلووات) & نقض وضعیت شارژ (درصد) & $P_{\mathrm{c}}<0$ (درصد) & $P_{\mathrm{d}}<0$ (درصد) & هم‌زمانی (درصد) \\
    \midrule
    مرجع پویا & $\approx 0$ & $0$ & $0$ & $0$ & $0$ \\
    بدون باتری & $\approx 0$ & $0$ & $0$ & $0$ & $0$ \\
    داده‌محور خام & $0{.}0166\pm 0{.}0021$ & $11{.}4\pm 2{.}8$ & $41{.}9\pm 7{.}0$ & $31{.}0\pm 9{.}0$ & $35{.}4\pm 3{.}7$ \\
    فیزیک‌تعبیه‌شده & $\approx 0$ & $0$ & $0$ & $0$ & $0$ \\
    \bottomrule
  \end{tabular}
\end{table}

شکل~\ref{fig:power-residual} قدر مطلق باقیمانده توازن را در هفته نخست آزمون نشان می‌دهد. شکل~\ref{fig:first-test-day} نخستین روز آزمون را جداگانه می‌آورد؛ این روز از مجموعه آزمون است، نه از آموزش یا اعتبارسنجی.

\begin{figure}[H]
  \centering
  \includegraphics[width=0.92\textwidth]{fig-first-test-day.png}
  \caption{بار و فتوولتائیک مطالعه موردی و وضعیت شارژ کنترل‌کننده‌ها در نخستین روز آزمون}
  \label{fig:first-test-day}
\end{figure}

\begin{figure}[H]
  \centering
  \includegraphics[width=0.92\textwidth]{fig-power-residual.png}
  \caption{قدر مطلق باقیمانده توازن توان در هفته نخست آزمون}
  \label{fig:power-residual}
\end{figure}

\section{هزینه اجرایی}
هزینه کل افق آزمون ۶۵۹ ساعته در جدول~\ref{tab:operating-cost} آمده است. هزینه مدل داده‌محور فقط پس از نگاشت ایمنی محاسبه شده تا نقض قید خام پاداش اقتصادی نگیرد. شبکه فیزیک‌تعبیه‌شده با هزینه $42{.}418\pm 0{.}511$ حدود $6{.}2$ درصد از حالت بدون باتری ارزان‌تر و به‌طور میانگین $2{.}8$ درصد از مرجع برنامه‌ریزی پویا گران‌تر است. هزینه آن با مدل داده‌محور ترمیم‌شده ($42{.}234\pm 0{.}318$) قابل مقایسه است و میانگین ترمیم‌شده اندکی کمتر است. با همپوشانی انحراف معیار پنج بذر، این اختلاف از نظر توصیفی قاطع نیست و آزمون فرض جداگانه‌ای گزارش نشده است. پس امکان‌پذیری تضمین‌شده به معنی رسیدن خودکار به برنامه بهینه افق کامل نیست. شکل~\ref{fig:daily-cost} میانگین هزینه روزانه برنامه‌های قابل اجرا را نشان می‌دهد.

زمان اجرا در محیط بایگانی‌شده آزمایش اندازه‌گیری شد: پردازنده شصت‌وچهاربیتی که مدل دقیق آن در آن محیط آشکار نبود، مفسر پایتون $3{.}11{.}2$ و کتابخانه محاسباتی نام‌پای $2{.}4{.}6$. حل یک‌باره مرجع برنامه‌ریزی پویا برای افق کامل ۳۲۸۹ گام $0{.}540$ ثانیه طول کشید. زمان یک گام کنترل‌کننده، شامل گذر پیش‌رو در شبکه $32\times 32$، نگاشت تانژانت هیپربولیک و لایه تحلیلی، پس از ۱۰۰۰ فراخوانی گرم‌سازی و در ۱۰ تکرار ۱۰۰۰۰تایی برابر $21{.}468\pm 0{.}402$ میکروثانیه بود. این اعداد با فرمان \lr{python code/benchmark\_runtime.py} بازتولید می‌شوند. برای سخت‌افزار دیگر باید نوع پردازنده یا شتاب‌دهنده، میانگین، انحراف معیار و شمار فراخوانی‌ها جداگانه ثبت شود.

\begin{table}[H]
  \centering
  \caption{هزینه اجرایی افق آزمون؛ واحد پولی مفروض}
  \label{tab:operating-cost}
  \begin{tabular}{lcc}
    \toprule
    مدل & هزینه کل آزمون & فاصله از مرجع پویا (درصد) \\
    \midrule
    مرجع پویا & $41{.}252$ & $0$ \\
    بدون باتری & $45{.}233$ & $9{.}7$ \\
    داده‌محور با نگاشت ایمنی پس‌پردازشی & $42{.}234\pm 0{.}318$ & $2{.}4\pm 0{.}8$ \\
    فیزیک‌تعبیه‌شده & $42{.}418\pm 0{.}511$ & $2{.}8\pm 1{.}2$ \\
    \bottomrule
  \end{tabular}
\end{table}

\begin{figure}[H]
  \centering
  \includegraphics[width=0.78\textwidth]{fig-daily-cost.png}
  \caption{میانگین هزینه روزانه برنامه‌های قابل اجرا روی افق آزمون}
  \label{fig:daily-cost}
\end{figure}

\section{تفسیر و محدودیت‌ها}
نتیجه اصلی فقط کاهش خطای میانگین نیست. شبکه فیزیک‌تعبیه‌شده با اندکی افزایش خطای تقلید نسبت به مدل آزاد، برنامه‌هایی می‌سازد که از نظر قیود تعریف‌شده همواره قابل اجرا هستند. مدل آزاد برای استفاده عملی به تعمیر پس‌پردازشی نیاز دارد. این موازنه با ادبیات قید سخت سازگار است: محدود کردن فضای خروجی آزادی تقریب را کم می‌کند، اما اعتماد فیزیکی را بالا می‌برد \pcite{lu2021hardconstraints,chen2024hardlinear}.

محدودیت‌ها نیز روشن‌اند. داده ترکیبی است و تابش فتوولتائیک ساعتی در همان خانه اندازه‌گیری نشده است. ورودی‌های یک و سه ساعت آینده و افق مرجع برنامه‌ریزی پویا با آینده‌نگری کامل‌اند. تبادل شبکه نامحدود فرض شده؛ اگر سقف اتصال باشد، لایه سخت باید فرمان باتری را با بازه مجاز $P_{\mathrm{G}}$ سازگار کند. مدل باتری دما، پیری، وابستگی بازده به توان و دینامیک مبدل را ندارد. نگاه‌به‌جلو در چند نمونه پایانی هر افراز می‌تواند سیگنال برون‌زاد افراز بعدی را ببیند؛ این نشت مرزی محدود است و برچسب آزمون در آموزش نیست. در آموزش $s$ از مسیر مرجع است و در آزمون بازگشتی از خود مدل؛ این ناهماهنگی توزیع هم در شکاف هزینه نقش دارد. بنابراین این گزارش برتری صنعتی یا عملکرد با خطای پیش‌بینی را نشان نمی‌دهد.

آنچه در ادامه می‌آید فقط پروتکل پیشنهادی کارهای آینده است، نه نتیجه آزمایش‌شده. در یک \textbf{جاروب حساسیت} هر بار فقط یکی از $P_{\max}$، $E_{\max}$، بازده‌ها، $\kappa$ و فاصله قیمت خرید و فروش، روی شبکه‌ای از پیش ثبت‌شده، عوض شود؛ برای هر حالت مرجع و پنج بذر دوباره اجرا شوند و هزینه، فاصله از برنامه‌ریزی پویا و نقض‌ها گزارش گردند. سپس ورودی‌های نگاه‌به‌جلو $t+1$ و $t+3$ جداگانه و هم‌زمان حذف شوند؛ پیش‌بینی‌ها با خطاهای کالیبره‌شده جایگزین گردند؛ و روزهای خارج از توزیع، برحسب فصل یا صدک بار، کنار گذاشته و فقط برای آزمون به کار روند. منحنی آموزش باید زیان آموزش و اعتبارسنجی هر بذر و نقطه توقف را نشان دهد. \textbf{تحلیل رژیمی} نیز، پیش از دیدن نتیجه، ساعت‌های نزدیک $s_{\min}$ و ساعت‌های تغییر دوره تعرفه را تعریف کند و میانگین قدر مطلق خطا، هزینه و نرخ نقض را جداگانه گزارش نماید.
